\section{Kepler bordered determinant and period-one conversion}\label{bpI:p6}

This part proves the frozen I-KEPLER (Theorem~\ref{bpI:p0:kepler}) completely, and supplies the bridge from the Kepler
Hamiltonian in the conventional variables ($2\pi$-periodic angles, actions $(L,G)$) to the frozen period-one
conventions of I-G1 and I-HEAD: the transformed Hamiltonian is again of Kepler type, with the constants
$\kappa'=\kappa(2\pi)^2$, $\Omega'=\Omega/(2\pi)$ (finding I06 of the Part-0 review: $\kappa,\Omega$ are kept general),
and its bordered Hessian \Lean{borderedHessian} is exactly the frozen matrix \Lean{keplerBorder} of the converted constants.
Part~\ref{bpI:p7} uses the result for the nondegeneracy hypothesis of I-HEAD; track~P cites
Lemma~\ref{bpI:p6:conversion} for the action normalisation.

\subsection{The determinant}\label{bpI:p6:sec-det}

\begin{lemma}[Determinant of the Kepler bordered matrix]\label{bpI:p6:det-formula}
For all $\kappa,\Omega\in\R$ and $L>0$ (no sign condition on $\kappa,\Omega$),
\[
\det\mathsf K_{\kappa,\Omega}(L)=\frac{3\kappa\Omega^2}{L^4}.
\]
\source{Chi05 p.~17, display after (3.43), for $\kappa=\Omega=1$ ($\det=3/L^4$); the general case is the same computation (ours).}
\end{lemma}
\begin{proof}
\proofstep{Expansion}\label{stepI:p6:expand}
Write $\mathsf K=(k_{ij})$ with $k_{11}=-3\kappa/L^4$, $k_{12}=0$, $k_{13}=\kappa/L^3$, $k_{21}=0$, $k_{22}=0$, $k_{23}=-\Omega$,
$k_{31}=\kappa/L^3$, $k_{32}=-\Omega$, $k_{33}=0$ (Definition~\ref{bpI:p0:keplerborder}; the denominators are
nonzero because $L>0$, the carrier of \Lean{keplerBorder} being $\{L:\R\mid0<L\}$). The Leibniz (Sarrus) expansion of a
$3\times3$ determinant (\Lean{Matrix.det\_fin\_three}) is
\[
\det\mathsf K=k_{11}k_{22}k_{33}+k_{12}k_{23}k_{31}+k_{13}k_{21}k_{32}-k_{13}k_{22}k_{31}-k_{12}k_{21}k_{33}-k_{11}k_{23}k_{32}.
\]
\proofstep{Evaluation}\label{stepI:p6:eval}
Since $k_{12}=k_{21}=k_{22}=k_{33}=0$, the first five products vanish (each contains one of these factors), and
$-k_{11}k_{23}k_{32}=-\bigl(-3\kappa/L^4\bigr)(-\Omega)(-\Omega)=3\kappa\Omega^2/L^4$.
(Equivalently: expand along the first column or second row; the cofactor of the $(1,3)$-entry has the vanishing minor
$\det\bigl(\begin{smallmatrix}0&0\\ \kappa/L^3&-\Omega\end{smallmatrix}\bigr)=0$, so only the $(1,1)$-entry contributes,
with minor $\det\bigl(\begin{smallmatrix}0&-\Omega\\-\Omega&0\end{smallmatrix}\bigr)=-\Omega^2$.)
\end{proof}

\begin{theorem}[I-KEPLER, proof]\label{bpI:p6:kepler}
For $\kappa,\Omega>0$ and $L>0$: $\det\mathsf K_{\kappa,\Omega}(L)=3\kappa\Omega^2/L^4$ and $\det\mathsf K_{\kappa,\Omega}(L)>0$.
This is Theorem~\ref{bpI:p0:kepler}, with the same quantifiers: $\kappa,\Omega$ first, then $L$ in the subtype; no invertibility of any Hessian is assumed.
\frozen{KAM.Iso.kepler\_bordered\_det}
\source{Chi05 p.~17; general $\kappa,\Omega$ ours.}
\end{theorem}
\begin{proof}
\proofstep{Value}\label{stepI:p6:k-value} Lemma~\ref{bpI:p6:det-formula}.

\proofstep{Positivity}\label{stepI:p6:k-pos} $\kappa>0$, $\Omega^2>0$ (as $\Omega\ne0$) and $L^4>0$ (as $L>0$), so the quotient of the positive numerator
$3\kappa\Omega^2$ by $L^4$ is positive (\Lean{div\_pos}, \Lean{pow\_pos}).
\end{proof}

\begin{remark}[Lean totality]\label{bpI:p6:totality}
In the statement of Theorem~\ref{bpI:p6:kepler} the only divisions are by powers of $L=L.\mathrm{val}>0$; there is no \Lean{fderiv},
inverse or square root. The hypotheses $\kappa,\Omega>0$ are used only for positivity; the determinant formula holds for all real $\kappa,\Omega$.
\end{remark}

%---------------------------------------------------------------------------------
\subsection{The Kepler Hamiltonian as a function on $\C^2$; its bordered Hessian}\label{bpI:p6:sec-border}

\begin{definition}[The Kepler--rotation Hamiltonian]\label{bpI:p6:kep-ham}
For $\kappa,\Omega\in\R$ let $h_{\kappa,\Omega}:\C^2\to\C$,
\[
h_{\kappa,\Omega}(I)=-\frac{\kappa}{2I_1^2}-\Omega I_2\qquad(I_1\ne0),
\]
(Lean division: the value at $I_1=0$ is junk and is never used; $h_{\kappa,\Omega}$ is holomorphic on $\{I_1\ne0\}$, which is open).
At real $I$ with $I_1\neq0$ its value is real. In the conventional variables $(L,G)$ the Kepler Hamiltonian with the rotating-frame
term (Chi05 (3.43) at $\varepsilon=0$, $\kappa=\Omega=1$) is $h_{1,1}(L,G)=-1/(2L^2)-G$.
\source{Chi05 (3.43), p.~16 ($\kappa=\Omega=1$); ours for general constants.}
\end{definition}

\begin{lemma}[Frequency map and bordered Hessian of $h_{\kappa,\Omega}$]\label{bpI:p6:kep-border}
Let $\kappa,\Omega\in\R$ and $I\in\R^2$ with $I_1>0$. Then
\[
\Lean{frequencyMap}\ h_{\kappa,\Omega}\ I=\Bigl(\frac{\kappa}{I_1^3},\,-\Omega\Bigr),
\qquad
D\bigl(\Lean{frequencyMap}\ h_{\kappa,\Omega}\bigr)(I)=\begin{pmatrix}-3\kappa/I_1^4&0\\0&0\end{pmatrix}
\]
(the real Fr\'echet derivative, $a$-th component applied to $e_b$ being the $(a,b)$ entry), and with the reindexing
$\epsilon:\mathrm{Fin}\,3\simeq\mathrm{Fin}\,2\oplus\mathrm{Unit}$, $0,1\mapsto\mathrm{inl}\,0,\mathrm{inl}\,1$, $2\mapsto\mathrm{inr}$,
\[
\Lean{borderedHessian}\ h_{\kappa,\Omega}\ I=\mathrm{reindex}_\epsilon\ \mathsf K_{\kappa,\Omega}(I_1),
\]
and in particular
\[
\det\Lean{borderedHessian}\ h_{\kappa,\Omega}\ I=\frac{3\kappa\Omega^2}{I_1^4}.
\]
\frozen{KAM.Iso.borderedHessian, KAM.frequencyMap, KAM.Iso.keplerBorder}
\source{Chi05 p.~17 ($\kappa=\Omega=1$); ours.}
\end{lemma}
\begin{proof}
\proofstep{Real part near a real point}\label{stepI:p6:kb-real}
Let $V=\{J\in\R^2:J_1>0\}$, open. For $J\in V$ the complex number $h_{\kappa,\Omega}(\iota J)$ is real and equals
$\phi(J):=-\kappa/(2J_1^2)-\Omega J_2$ ($\iota:\R^2\to\C^2$ the inclusion, $\Lean{Complex.ofReal}$ applied entrywise), so
$J\mapsto\operatorname{Re}h_{\kappa,\Omega}(\iota J)$ coincides with the smooth function $\phi$ on the neighbourhood $V$ of $I$.

\proofstep{Gradient}\label{stepI:p6:kb-grad} Derivatives at a point depend only on the germ (\Lean{Filter.EventuallyEq.fderiv\_eq}); so
$\Lean{frequencyMap}\,h\,J_0=\nabla\phi(J_0)$ for $J_0\in V$, with $\partial_1\phi=\kappa/J_1^3$, $\partial_2\phi=-\Omega$ (power rule).
Thus on the open set $V$ the map $\Lean{frequencyMap}\,h$ coincides with $g(J)=(\kappa J_1^{-3},-\Omega)$.

\proofstep{Derivative of the frequency map}\label{stepI:p6:kb-hess} Again by germs, $D(\Lean{frequencyMap}\,h)(I)=Dg(I)$, whose first row is
$(-3\kappa I_1^{-4},0)$ and whose second row is $0$.

\proofstep{Entries of the bordered matrix}\label{stepI:p6:kb-entries} Insert into Definition~\ref{bpI:p0:bordered}: entries $(a,b)$ with
$a,b\in\{0,1\}$ are those of $Dg(I)$; the border entries $(a,\mathrm{inr})$ and $(\mathrm{inr},b)$ are $g(I)_a$ and $g(I)_b$, i.e. $\kappa/I_1^3,-\Omega$;
the corner is $0$. Compare entrywise with Definition~\ref{bpI:p0:keplerborder} at $L=I_1$: the matrices agree (the matrix $\mathsf K$ is symmetric,
the bordered Hessian was not assumed to be).

\proofstep{Determinant}\label{stepI:p6:kb-det} The determinant is invariant under reindexing by a bijection of the index set (\Lean{Matrix.det\_reindex\_self}),
and Lemma~\ref{bpI:p6:det-formula} gives the value.
\end{proof}

\subsection{Conversion to period-one angles}\label{bpI:p6:sec-convert}

\begin{definition}[Period-one conversion map]\label{bpI:p6:conversion-map}
$\mathsf S:\R^2\times\R^2\to\R^2\times\R^2$, $\mathsf S(\theta,I)=\bigl(2\pi\theta,\ I/(2\pi)\bigr)$, so that
$(\ell,g)=2\pi\theta$ and $(L,G)=I/(2\pi)$, i.e. $\theta=(\ell,g)/(2\pi)$ and $I=2\pi(L,G)$; the same formula on $\C^2\times\C^2$.
It is a linear isomorphism, and it maps $\Z^2$-translations of $\theta$ to $2\pi\Z^2$-translations of $(\ell,g)$.
\source{ours; frozen convention of \Kref{bp:p0:conventions}.}
\end{definition}

\begin{lemma}[Conversion]\label{bpI:p6:conversion}
Let $\kappa,\Omega\in\R$, and $\kappa'=\kappa(2\pi)^2$, $\Omega'=\Omega/(2\pi)$.
\begin{enumerate}[label=(\alph*)]
\item \emph{Symplectic.} $\mathsf S^*\bigl(d\ell_1\wedge dL_1+d\ell_2\wedge dL_2\bigr)=d\theta_1\wedge dI_1+d\theta_2\wedge dI_2$, i.e. with $\Omega_0=\begin{pmatrix}0&\mathbf 1\\-\mathbf 1&0\end{pmatrix}$
  (the matrix of the standard form in the order angles, actions): $D\mathsf S^{\top}\Omega_0D\mathsf S=\Omega_0$ and $D\mathsf S\,\Omega_0\,D\mathsf S^{\top}=\Omega_0$.
  Consequently, for $H$ differentiable, $D\mathsf S\,X_{H\circ\mathsf S}=X_H\circ\mathsf S$ (with $X_H=(\partial_IH,-\partial_\theta H)$),
  and a curve $w$ solves $\dot w=X_{H\circ\mathsf S}(w)$ iff $\mathsf Sw$ solves $\dot y=X_H(y)$.
\item \emph{Hamiltonian.} For all $I\in\C^2$ with $I_1\neq0$: $h_{\kappa,\Omega}\bigl(I/(2\pi)\bigr)=h_{\kappa',\Omega'}(I)$. Thus $\tilde h:=h_{\kappa,\Omega}\circ(I\mapsto I/2\pi)$ is the Kepler--rotation Hamiltonian
  of the converted constants, holomorphic on $\{I_1\ne0\}$, real at real $I$. For $\kappa=\Omega=1$: $\kappa'=(2\pi)^2>0$, $\Omega'=1/(2\pi)>0$.
\item \emph{Frequencies.} For real $I$ with $I_1>0$, $\nabla\tilde h(I)=\bigl(\kappa'/I_1^3,-\Omega'\bigr)=\frac1{2\pi}\nabla h_{\kappa,\Omega}\bigl(I/2\pi\bigr)$: the frequency per unit of the period-one angle is the conventional angular velocity divided by $2\pi$.
\item \emph{Bordered Hessian.} For real $I$ with $I_1>0$, with $\mathsf D=\mathrm{diag}\bigl(\tfrac1{2\pi},\tfrac1{2\pi},1\bigr)$,
  \[
  \begin{aligned}
  \mathsf K_{\kappa',\Omega'}(I_1)&=\mathsf D\,\mathsf K_{\kappa,\Omega}\bigl(I_1/2\pi\bigr)\,\mathsf D,\qquad
  \Lean{borderedHessian}\,\tilde h\,I=\mathrm{reindex}_\epsilon\,\mathsf K_{\kappa',\Omega'}(I_1),\\
  \det\Lean{borderedHessian}\,\tilde h\,I&=\frac{3\kappa'\Omega'^2}{I_1^4}=\frac{3\kappa\Omega^2}{I_1^4}.
  \end{aligned}
  \]
  In particular for $\kappa=\Omega=1$: $\det\Lean{borderedHessian}\,\tilde h\,I=3/I_1^4>0$.
\end{enumerate}
\source{ours; the conversion by $2\pi$ is forced by the period-one convention of \Kref{bp:p0:conventions} (Chi05 uses $2\pi$-periodic angles).
Finding I06 of the Part-0 review.}
\end{lemma}
\begin{proof}
\proofstep{Symplectic}\label{stepI:p6:cv-sympl} $D\mathsf S=\mathrm{diag}(2\pi,2\pi,\tfrac1{2\pi},\tfrac1{2\pi})$ in the order $(\theta_1,\theta_2,I_1,I_2)$. A block computation gives
$D\mathsf S^{\top}\Omega_0D\mathsf S=\begin{pmatrix}0&2\pi\cdot\frac1{2\pi}\mathbf 1\\-\frac1{2\pi}\cdot2\pi\mathbf 1&0\end{pmatrix}=\Omega_0$, and likewise for $D\mathsf S\Omega_0D\mathsf S^{\top}$ (the matrix is diagonal, hence equal to its transpose).
Since $X_H=\Omega_0\nabla H$ and $\nabla(H\circ\mathsf S)=D\mathsf S^{\top}\nabla H\circ\mathsf S$:
$D\mathsf S\,X_{H\circ\mathsf S}=D\mathsf S\,\Omega_0D\mathsf S^{\top}\nabla H\circ\mathsf S=\Omega_0\nabla H\circ\mathsf S$. For a curve, the chain rule gives
$\frac{d}{dt}(\mathsf Sw)=D\mathsf S\dot w$.

\proofstep{Hamiltonian}\label{stepI:p6:cv-ham} For $I_1\ne0$: $-\kappa/(2(I_1/2\pi)^2)=-\kappa(2\pi)^2/(2I_1^2)$ and $\Omega\,I_2/(2\pi)=\Omega'I_2$ (field arithmetic).

\proofstep{Frequencies and bordered Hessian}\label{stepI:p6:cv-border} Part (c) is Lemma~\ref{bpI:p6:kep-border} applied to $(\kappa',\Omega')$ together with (b) and the chain rule for $I\mapsto I/2\pi$.
For the matrix identity: $(\mathsf D\mathsf K\mathsf D)_{ij}=\mathsf D_{ii}\mathsf K_{ij}\mathsf D_{jj}$ with $L=I_1/2\pi$:
$(1,1)$: $-3\kappa/(L^4(2\pi)^2)=-3\kappa(2\pi)^2/I_1^4$; $(1,3)$: $\kappa/(L^3\,2\pi)=\kappa(2\pi)^2/I_1^3$; $(2,3)$: $-\Omega/(2\pi)=-\Omega'$; the other entries are $0$. These are the entries of $\mathsf K_{\kappa',\Omega'}(I_1)$.
The determinant: $3\kappa'\Omega'^2/I_1^4=3\kappa(2\pi)^2\Omega^2(2\pi)^{-2}/I_1^4=3\kappa\Omega^2/I_1^4$; also $\det(\mathsf D\mathsf K\mathsf D)=(2\pi)^{-4}\cdot3\kappa\Omega^2(2\pi)^4/I_1^4$, the same value.
\end{proof}

\begin{corollary}[Nondegeneracy in period-one variables]\label{bpI:p6:nondeg}
With $\kappa=\Omega=1$, $\tilde h=h_{(2\pi)^2,1/(2\pi)}$ is holomorphic on $\{I\in\C^2:I_1\ne0\}$, real at real $I$, and $\det\Lean{borderedHessian}\,\tilde h\,I=3/I_1^4>0$ for every real $I$ with $I_1>0$. If $B\subset\{I\in\R^2:I_1>0\}$ and $0<r_0\le\tfrac12\inf_BI_1$, then $\tilde h$ is holomorphic on $B_{r_0}$, every real $I\in B_{r_0}$ has $I_1>0$, and the nondegeneracy hypothesis of Theorem~\ref{bpI:p0:head} ($\det\Lean{borderedHessian}\,\tilde h\,I\ne0$ for every real $I\in B_{r_0}$) holds.
\end{corollary}
\begin{proof}
\proofstep{Domain}\label{stepI:p6:nd-domain} If $I\in B_{r_0}$ then $\norm{I-J}<r_0$ for some $J\in B$, so $\abs{\operatorname{Re}I_1-J_1}<r_0\le\tfrac12J_1$ and $\operatorname{Re}I_1>\tfrac12J_1>0$; in particular $I_1\ne0$, and $I_1>0$ for real $I$.
\proofstep{Conclusion}\label{stepI:p6:nd-concl} Lemma~\ref{bpI:p6:conversion}(b),(d) and Lemma~\ref{bpI:p6:kep-border}.
\end{proof}

\begin{remark}[What is used where]\label{bpI:p6:uses}
Theorem~\ref{bpI:p6:kepler} is the whole of I-KEPLER (no consumer needs more). Corollary~\ref{bpI:p6:nondeg} is the nondegeneracy hypothesis \texttt{hiso} of I-HEAD for the Kepler part of the restricted problem (Part~\ref{bpI:p7}, Proposition~\ref{bpI:p7:nondeg}); Lemma~\ref{bpI:p6:conversion}(a) is the symplectic bookkeeping for the passage between the Delaunay chart of Part~\ref{bpI:p7} (conventional $2\pi$ variables) and I-HEAD (period one).
Track~P cites Lemma~\ref{bpI:p6:conversion} for the action normalisation. No external (Mathlib-level) fact beyond field arithmetic, the determinant of a $3\times3$ matrix, the power rule and the germ locality of the Fr\'echet derivative is used.
\end{remark}
